\documentclass[11pt]{article}

\usepackage[margin=1.02in]{geometry}
\usepackage{amsmath,amssymb,amsthm,mathtools}
\usepackage{booktabs}
\usepackage{enumitem}
\usepackage{microtype}
\usepackage{xcolor}
\usepackage[hidelinks]{hyperref}
\usepackage{fancyhdr}

\definecolor{deepblue}{HTML}{173B57}
\hypersetup{
  colorlinks=true,
  linkcolor=deepblue,
  citecolor=deepblue,
  urlcolor=deepblue,
  pdftitle={Maximal violation without maximal global randomness in a cyclic Bell family},
  pdfauthor={Alec Kriebel}
}

\pagestyle{fancy}
\fancyhf{}
\fancyhead[L]{\small\textsc{Cyclic Bell maximizers}}
\fancyhead[R]{\small\textsc{A. Kriebel}}
\fancyfoot[C]{\thepage}
\setlength{\headheight}{14pt}

\newtheorem{theorem}{Theorem}[section]
\newtheorem{lemma}[theorem]{Lemma}
\newtheorem{proposition}[theorem]{Proposition}
\newtheorem{corollary}[theorem]{Corollary}
\theoremstyle{remark}
\newtheorem{remark}[theorem]{Remark}

\newcommand{\Id}{\mathbf 1}
\newcommand{\cI}{\mathcal I}
\newcommand{\cH}{\mathcal H}
\newcommand{\Reop}{\operatorname{Re}}
\newcommand{\Tr}{\operatorname{Tr}}
\newcommand{\supp}{\operatorname{supp}}
\newcommand{\e}{\mathrm e}
\newcommand{\bbC}{\mathbb C}
\newcommand{\absop}[1]{\lvert #1\rvert}

\title{\vspace{-1.4cm}\textbf{Maximal violation without maximal global randomness}\\[0.25em]
\textbf{in a cyclic Bell family}\\[0.45em]
\large Exact counterexamples in every dimension \(d\geq4\)}
\author{Alec Kriebel\\
\small with heavy assistance from ChatGPT 5.6 Sol}
\date{26 July 2026}

\begin{document}
\maketitle

\begin{abstract}
Perito, D'Avino, Jung, Mironowicz, Acin, and Augusiak conjectured that
maximal violation of an augmented cyclic Bell functional certifies
\(2\log_2d\) bits of global randomness.  We disprove the
maximal-randomness assertion for every \(d\geq4\).  For each such \(d\), a
sparse weighted-shift strategy on
\(\bbC^d\otimes\bbC^d\), with no adversarial side information, attains the
exact Bell maximum
\[
  2\csc\!\left(\frac{\pi}{2d}\right)+1
\]
while the designated joint output distribution is nonuniform.  Quantitatively,
\[
  G(AB\mid1,d,E)
  \geq\frac1{d^2}
  +\frac{2\sin(\pi/d)\sin(3\pi/d)}{d^2(d-1)}
  >\frac1{d^2}.
\]
Both one-party marginals remain uniform, so the obstruction concerns global
joint randomness rather than a local-output bias.  At \(d=4\), a compact
certificate over \(\mathbb Q(\zeta_{16})\) gives probabilities \(1/32\) and
\(3/32\), hence \(G=3/32>1/16\).  We reproduce the exact Bell upper
certificate, prove the infinite counterexample family from first principles,
and provide an exact-arithmetic verifier for the finite witness.
\end{abstract}

\medskip
\noindent\textbf{Status.}
This is an unrefereed, AI-assisted research note.  Its proof and exact
certificate are public for audit but have not received independent expert
review.

\section{Definitions and result}

Fix \(d\geq2\) and \(\omega=\e^{2\pi i/d}\).  Alice has order-\(d\)
unitaries \(A_0,A_1\), Bob has order-\(d\) unitaries
\(B_0,\ldots,B_d\), and
\begin{align}
  \cI_d
  &=
  \frac12\sum_{x=0}^{1}\sum_{y=0}^{d-1}
  \left(
    \omega^{xy}A_x\otimes B_y+
    \omega^{-xy}A_x^\dagger\otimes B_y^\dagger
  \right),\label{eq:I}\\
  \overline{\cI}_d
  &=
  \cI_d+
  \frac12\left(
    A_0\otimes B_d+A_0^\dagger\otimes B_d^\dagger
  \right).\label{eq:Ibar}
\end{align}
These are Eqs.~(11) and (10), respectively, of
Ref.~\cite{PeritoEtAl2026}.  Put
\begin{equation}
  M_d=2\csc\!\left(\frac{\pi}{2d}\right).
  \label{eq:Md}
\end{equation}
The exact operator bound
\(\cI_d\leq M_d\Id\), proved in a companion note
\cite{KriebelBound2026}, gives
\begin{equation}
  \overline{\cI}_d\leq(M_d+1)\Id.
  \label{eq:bar-bound}
\end{equation}
Section~\ref{sec:upper} recalls a self-contained proof.

For the output settings \(x=1,y=d\), let
\(\sigma_E^{ab}\) be Eve's subnormalized conditional states.  Their guessing
probability is
\[
  G(AB\mid1,d,E)
  =
  \max_{\{Q_{ab}\}}\sum_{a,b}\Tr(Q_{ab}\sigma_E^{ab}).
\]
Maximal private global randomness would require
\(\sigma_E^{ab}=\rho_E/d^2\), and hence \(G=1/d^2\).

\begin{theorem}[Failure of maximal global randomness]\label{thm:main}
For every integer \(d\geq4\), there is an explicit admissible strategy on
\(\bbC^d\otimes\bbC^d\), with \(\cH_E=\bbC\), such that
\[
  \langle\overline{\cI}_d\rangle=M_d+1
\]
but
\begin{equation}
  G(AB\mid1,d,E)
  \geq\frac1{d^2}
  +\frac{2\sin(\pi/d)\sin(3\pi/d)}{d^2(d-1)}
  >\frac1{d^2}.
  \label{eq:guess-gap}
\end{equation}
For the same strategy, each Alice and Bob marginal is uniform.  Thus maximal
violation does not force maximal observed joint uniformity, and therefore
cannot force maximal private global randomness.
\end{theorem}

Under the operator definitions \eqref{eq:I}--\eqref{eq:Ibar}, the maximum in
Theorem~\ref{thm:main} is fixed by \eqref{eq:bar-bound} and an attaining
strategy.  Equation~(17) of Ref.~\cite{PeritoEtAl2026} likewise gives
\(M_d\) for the unaugmented strategy, and its \(d=3\) discussion gives the
augmented value \(5=M_3+1\).  The sentence preceding Conjecture~2 and the
conjecture itself print an additional factor \(d\) in the first term.  We
therefore treat the displayed operators and Eq.~(17) as authoritative.  The
counterexample addresses the conjecture's substantive assertion that
\emph{maximal violation} of the first augmented family certifies
\(2\log_2d\) random bits.

\section{The exact Bell upper certificate}\label{sec:upper}

We reproduce the argument because global optimality is essential to the
counterexample.  Write
\[
  C_y=A_0+\omega^yA_1,\qquad U=A_0^\dagger A_1,\qquad
  F_d(U)=\sum_{y=0}^{d-1}\absop{\Id+\omega^yU}.
\]

\begin{lemma}[Polar positive-factor identity]\label{lem:polar}
Let \(C=V\absop C\) be the polar decomposition of a finite-dimensional
operator; \(V\) may be a partial isometry.  If \(B\) is unitary, then
\[
  P=\absop{C^\dagger}^{1/2}\otimes\Id
    -V\absop C^{1/2}\otimes B
\]
satisfies
\begin{equation}
  \frac{\absop C+\absop{C^\dagger}}2\otimes\Id
  -\Reop(C\otimes B)
  =\frac12P^\dagger P.
  \label{eq:polar}
\end{equation}
\end{lemma}

\begin{proof}
The initial projection \(V^\dagger V\) is the support projection of
\(\absop C\), while polar functional calculus gives
\(\absop{C^\dagger}^{1/2}V=V\absop C^{1/2}\).  Consequently,
\[
\begin{aligned}
 P^\dagger P
 &=
 \absop{C^\dagger}\otimes\Id+
 \absop C^{1/2}V^\dagger V\absop C^{1/2}\otimes B^\dagger B\\
 &\quad-C\otimes B-C^\dagger\otimes B^\dagger\\
 &=
 (\absop C+\absop{C^\dagger})\otimes\Id
 -C\otimes B-C^\dagger\otimes B^\dagger.
\end{aligned}
\]
This also shows explicitly why no invertibility assumption is needed.
\end{proof}

\begin{lemma}[Sharp scalar bound]\label{lem:scalar}
For \(\lvert z\rvert=1\),
\begin{equation}
  \sum_{y=0}^{d-1}\lvert1+\omega^yz\rvert\leq M_d,
  \label{eq:scalar}
\end{equation}
with equality exactly when
\begin{equation}
  z^d=(-1)^{d-1}.
  \label{eq:equality-roots}
\end{equation}
\end{lemma}

\begin{proof}
Write \(z=\e^{2is}\), \(h=\pi/d\), and \(a=h/2\).  Half the left side is
\[
  G_d(s)=\sum_{y=0}^{d-1}\lvert\cos(s+yh)\rvert,
\]
which has period \(h\).

If \(d=2m+1\), take \(-a\leq s\leq a\) and reindex by
\(k=-m,\ldots,m\).  Every \(s+kh\) lies in
\([-\pi/2,\pi/2]\), since \(a+mh=\pi/2\), and the finite cosine sum gives
\[
  G_d(s)=\sum_{k=-m}^{m}\cos(s+kh)=\csc(a)\cos s.
\]
Its maximum occurs exactly at \(s\equiv0\pmod h\), equivalently \(z^d=1\).

If \(d=2m\), take \(0\leq s\leq h\).  The first \(m\) cosines are
nonnegative and the last \(m\) nonpositive; endpoints cause no ambiguity
because the boundary cosine vanishes.  Two finite cosine sums give
\[
  G_d(s)
  =
  \sum_{y=0}^{m-1}\cos(s+yh)
  -\sum_{y=m}^{2m-1}\cos(s+yh)
  =\csc(a)\cos(s-a).
\]
Its maximum occurs exactly at \(s\equiv a\pmod h\), equivalently \(z^d=-1\).
Multiplying by two proves both claims.
\end{proof}

Because \(\Id+\omega^yU\) is normal,
\[
  \absop{C_y}=\absop{\Id+\omega^yU},\qquad
  \absop{C_y^\dagger}
  =A_0\absop{\Id+\omega^yU}A_0^\dagger.
\]
Functional calculus and Lemma~\ref{lem:scalar} give
\(F_d(U)\leq M_d\Id\).  Summing Lemma~\ref{lem:polar} therefore yields the
exact gap decomposition
\begin{equation}
\boxed{
\begin{aligned}
 M_d\Id-\cI_d
 ={}&\frac12\sum_{y=0}^{d-1}P_y^\dagger P_y\\
 &+\frac12\left[
 M_d\Id-F_d(U)
 +A_0\bigl(M_d\Id-F_d(U)\bigr)A_0^\dagger
 \right]\otimes\Id.
\end{aligned}}
\label{eq:gap}
\end{equation}
Every term is positive.  This proves
\(\cI_d\leq M_d\Id\) for arbitrary unitaries, even without the order
relations, and \eqref{eq:bar-bound} follows from
\(\Reop(A_0\otimes B_d)\leq\Id\).

\section{Weighted-shift maximizers}\label{sec:family}

Let
\[
  X\lvert j\rangle=\lvert j+1\bmod d\rangle,\qquad
  \delta_d=\begin{cases}0,&d\text{ odd},\\1,&d\text{ even},\end{cases}
\]
and enumerate the equality roots from Lemma~\ref{lem:scalar} by
\begin{equation}
  z_k=\exp\!\left(\frac{\pi i(2k+\delta_d)}d\right),
  \qquad k=0,\ldots,d-1.
  \label{eq:zk}
\end{equation}
For any ordering
\(\kappa=(\kappa_0,\ldots,\kappa_{d-1})\) of these roots, define
\begin{align*}
 Z_\kappa&=\operatorname{diag}
 (z_{\kappa_0},\ldots,z_{\kappa_{d-1}}),\\
 s_{y,k}&=\frac{1+\omega^yz_k}{\lvert1+\omega^yz_k\rvert},\\
 S_y^\kappa&=\operatorname{diag}
 (s_{y,\kappa_0},\ldots,s_{y,\kappa_{d-1}}),\\
 H_y^\kappa&=\operatorname{diag}
 (\lvert1+\omega^yz_{\kappa_0}\rvert,\ldots,
  \lvert1+\omega^yz_{\kappa_{d-1}}\rvert).
\end{align*}
The denominators never vanish because \(-1\) is not a root of
\(t^d-(-1)^{d-1}\).

\begin{proposition}[Every root ordering gives a maximizer]
\label{prop:maximizer}
On \(\bbC^d\otimes\bbC^d\), set
\begin{equation}
\begin{aligned}
 A_0&=X,& A_1&=XZ_\kappa,\\
 V_y&=XS_y^\kappa,& B_y&=\overline{V_y}\quad(0\leq y<d),\\
 && B_d&=X,
\end{aligned}
\label{eq:strategy}
\end{equation}
and use
\[
  \lvert\Phi_d\rangle
  =\frac1{\sqrt d}\sum_{j=0}^{d-1}\lvert j,j\rangle.
\]
All observables are unitary, all have \(d\)-th power \(\Id\), and
\[
  \langle\Phi_d\vert\overline{\cI}_d\vert\Phi_d\rangle=M_d+1.
\]
\end{proposition}

\begin{proof}
The \(z_k\) are all roots of \(t^d-(-1)^{d-1}\).  Comparing the constant
term and evaluating the polynomial at \(-1\) give
\begin{equation}
  \prod_kz_k=1,\qquad
  \prod_k(1+\omega^yz_k)=2.
  \label{eq:products}
\end{equation}
Multiplication by \(\omega^y\) merely permutes the roots.  The second product
is real and positive, so
\[
  \prod_ks_{y,k}=1.
\]
For arbitrary weights \(t_j\),
\begin{equation}
  \left(X\operatorname{diag}(t_0,\ldots,t_{d-1})\right)^d
  =\left(\prod_jt_j\right)\Id.
  \label{eq:shift-power}
\end{equation}
Equations \eqref{eq:products}--\eqref{eq:shift-power} prove every order
relation in \eqref{eq:strategy}; unitarity is immediate.

Moreover,
\[
  A_0+\omega^yA_1
  =X(\Id+\omega^yZ_\kappa)
  =V_yH_y^\kappa,
\]
which is its polar decomposition.  The maximally entangled trace identity
\[
  \langle\Phi_d\vert R\otimes S\vert\Phi_d\rangle
  =\frac1d\Tr(RS^T)
\]
and \(B_y^T=V_y^\dagger\) give
\[
  \langle\cI_d\rangle
  =\frac1d\sum_{j,y}\lvert1+\omega^yz_{\kappa_j}\rvert
  =\frac1d\sum_jM_d=M_d.
\]
Finally, \(X\otimes X\) fixes \(\lvert\Phi_d\rangle\), so the added
correlator equals one.  The upper bound \eqref{eq:bar-bound} proves global
optimality.
\end{proof}

\section{The designated output distribution}

The eigenvectors of \(A_0=X\) and \(A_1=XZ_\kappa\) can be written
\[
  u_c=\frac1{\sqrt d}\sum_{j=0}^{d-1}\omega^{-cj}\lvert j\rangle,
  \qquad
  v_a=\frac1{\sqrt d}\sum_{j=0}^{d-1}q_j\omega^{-aj}\lvert j\rangle,
\]
where
\begin{equation}
  q_0=1,\qquad q_{j+1}=z_{\kappa_j}q_j.
  \label{eq:q}
\end{equation}
The sequence is cyclic because \(\prod_kz_k=1\).  Put
\[
  \widehat q_m=\sum_{j=0}^{d-1}q_j\omega^{mj}.
\]
The transpose of the outcome-\(b\) projector of \(B_d=X\) is the
outcome-\((-b)\) projector of \(A_0\).  Hence
\begin{equation}
  p(a,b\mid1,d)
  =\frac1d\lvert\langle u_{-b},v_a\rangle\rvert^2
  =\frac{\lvert\widehat q_{-(a+b)}\rvert^2}{d^3}.
  \label{eq:probability}
\end{equation}
Parseval gives
\(\sum_m\lvert\widehat q_m\rvert^2=d^2\).  As \(a+b\) traverses all
residues when either output is fixed, \eqref{eq:probability} immediately
implies
\begin{equation}
  \sum_bp(a,b\mid1,d)=\sum_ap(a,b\mid1,d)=\frac1d.
  \label{eq:marginals}
\end{equation}

For \(d\geq4\), choose
\begin{equation}
  \kappa=(0,1,\ldots,d-3,d-1,d-2),
  \label{eq:bad-order}
\end{equation}
the canonical ordering with its final two entries exchanged.  Consider the
second cyclic autocorrelation
\[
  C_2=\sum_jq_{j+2}\overline{q_j}
  =\sum_jz_{\kappa_j}z_{\kappa_{j+1}}.
\]
For the canonical ordering this sum vanishes.  Exchanging the last two
entries changes only the affected summands and gives
\begin{equation}
  C_2=(z_{d-1}-z_{d-2})(z_{d-3}-z_0),
  \qquad
  \lvert C_2\rvert
  =4\sin(\pi/d)\sin(3\pi/d)>0.
  \label{eq:C2}
\end{equation}
Thus the Fourier magnitudes cannot all equal \(d\), so the joint table is
not uniform.

For the quantitative gap, put
\[
  x_m=\lvert\widehat q_m\rvert^2-d,\qquad
  \Delta=\max_mx_m.
\]
Then \(\sum_mx_m=0\), and Fourier inversion at lag two gives
\[
  C_2=\frac1d\sum_mx_m\omega^{-2m}.
\]
Here the subtracted constant contributes zero because \(d\geq4\) and
\(\sum_m\omega^{-2m}=0\).
The nonzero mean-zero sequence \(x_m\) has at most \(d-1\) positive entries.
Therefore
\[
  \sum_m\lvert x_m\rvert
  =2\sum_{x_m>0}x_m
  \leq2(d-1)\Delta,
\]
and
\[
  \Delta\geq\frac{d\lvert C_2\rvert}{2(d-1)}.
\]
With trivial Eve, the optimal guessing strategy selects a most likely output
pair, so
\[
  G=\max_{a,b}p(a,b\mid1,d)=\frac{d+\Delta}{d^3}.
\]
Substitution of \eqref{eq:C2} proves \eqref{eq:guess-gap} and
Theorem~\ref{thm:main}.

\section{A sparse exact four-outcome certificate}\label{sec:d4}

Let \(\zeta=\e^{\pi i/8}\), so
\(\mathbb Q(\zeta)\cong\mathbb Q[x]/(x^8+1)\).  At \(d=4\), choose
\(\kappa=(0,1,3,2)\).  Proposition~\ref{prop:maximizer} becomes
\[
  A_0=X,\qquad
  A_1=X\operatorname{diag}(\zeta^2,\zeta^6,\zeta^{14},\zeta^{10}),
  \qquad B_4=X.
\]
The four polar unitaries
\(
V_y=X\operatorname{diag}(\zeta^{e_{y,0}},\ldots,\zeta^{e_{y,3}})
\)
are encoded by
\[
\begin{array}{c@{\qquad}c@{\qquad}c@{\qquad}c@{\qquad}c}
\toprule
y&e_{y,0}&e_{y,1}&e_{y,2}&e_{y,3}\\
\midrule
0&1&3&15&13\\
1&3&13&1&15\\
2&13&15&3&1\\
3&15&1&13&3\\
\bottomrule
\end{array}
\qquad B_y=\overline{V_y}.
\]
These sparse monomial formulas are exact matrices over
\(\mathbb Q(\zeta_{16})\); no numerical polar decomposition is part of the
certificate.

The sequence \eqref{eq:q} is
\[
  q=(1,\zeta^2,-1,\zeta^6).
\]
With \(\omega=i\), direct finite Fourier transformation gives
\[
  \widehat q=(i\sqrt2,\ 2+i\sqrt2,\ -i\sqrt2,\ 2-i\sqrt2),
  \qquad
  \bigl(\lvert\widehat q_m\rvert^2\bigr)_m=(2,6,2,6).
\]
Equation~\eqref{eq:probability} therefore gives
\begin{equation}
  p(a,b\mid1,4)=
  \begin{cases}
    1/32,&a+b\text{ even},\\
    3/32,&a+b\text{ odd}.
  \end{cases}
  \label{eq:d4-table}
\end{equation}
Consequently,
\[
  G(AB\mid1,4,E)=\frac3{32}>\frac1{16},
  \qquad
  H_{\min}(AB\mid E)=5-\log_2 3<4.
\]
The independent verifier reconstructs the observables from the displayed
exponents and checks, using exact rational arithmetic modulo \(x^8+1\),
all unitarity and fourth-order relations, all polar factorizations, the
original Bell expression, both probability calculations, uniform marginals,
and the strict guessing gap.

\section{Consequences, verification, and scope}

\begin{enumerate}[leftmargin=1.55em]
  \item The maximal-randomness assertion in Conjecture~2 of
  Ref.~\cite{PeritoEtAl2026}, interpreted for its displayed first augmented
  Bell family, is false for every \(d\geq4\).
  \item The obstruction needs neither a direct sum nor an Eve-held flag.
  It occurs on one \(d\times d\) block, on a maximally entangled state, while
  both local marginals remain uniform.
  \item The Bell maximum does not determine a unique target behavior in these
  dimensions.  In particular, the first augmented family cannot self-test a
  unique maximizing correlation.
  \item No universal near-maximal bound converging to \(1/d^2\) as the Bell
  deficit tends to zero can hold for this family when \(d\geq4\), because the
  exact maximizing face already contains nonuniform behaviors.
  \item The complete maximizing face, its worst guessing probability, and the
  maximal-violation randomness question for \(d=2,3\) remain unresolved here.
  The Appendix proves only a multiplicity restriction, not a classification.
\end{enumerate}

The source preprint's Appendix B.1 fixes the \emph{full canonical
probability distribution} before running its NPA guessing-probability
calculation.  That calculation can support randomness of the fixed canonical
behavior without excluding another behavior at the same Bell maximum.  The
strategy above supplies exactly such a behavior.  This note does not dispute
the separate protocols of Farkas, Mironowicz, and Augusiak
\cite{FarkasEtAl2026}, nor the distinct three-outcome robustness study of
D'Avino et al.~\cite{DavinoEtAl2026}.

The accompanying package contains the TeX source, the exact
\(\mathbb Q(\zeta_{16})\) verifier, a compact machine-readable certificate,
the all-dimensional constructor and symbolic family certificate, floating
point regressions for \(d=2,\ldots,12\), historical discovery code, claims and
assumptions ledgers, failed approaches, a research log, and SHA-256 hashes.
The finite programs support the proof but do not formally quantify over all
\(d\).

A targeted search on 26 July 2026 found no public equivalent of the
weighted-shift counterexample.  This does not establish absolute novelty:
unindexed, unpublished, or concurrent work may exist.  The appropriate
priority wording is therefore ``to the best of our knowledge.''  No
originating author or other researcher was contacted.

The proof architecture, derivations, manuscript, verification code, and
literature audit were developed with heavy assistance from ChatGPT 5.6 Sol
under Alec Kriebel's direction.  Alec Kriebel is the named human author but
cannot independently validate the mathematical claims.  Independent expert
checking remains necessary.

\appendix
\section{A support-level multiplicity restriction}\label{app:multiplicity}

This appendix records a structural consequence that is not needed for the
counterexample.

\begin{lemma}[Reflection-product rank bound]\label{lem:reflection}
Let \(T\) be a unitary on an \(n\)-dimensional space and let \(E\) be an
orthogonal projection of rank \(r\).  If
\[
  T^d=\Id,\qquad \bigl(T(\Id-2E)\bigr)^d=-\Id,
\]
then \(n\leq dr\).
\end{lemma}

\begin{proof}
Expanding the second product and using \(T^d=\Id\) expresses \(-\Id\) as a
product of \(d\) conjugates of the reflection \(\Id-2E\).  For arbitrary
operators \(R_1,\ldots,R_d\),
\[
  \operatorname{rank}(\Id-R_1\cdots R_d)
  \leq\sum_{j=1}^d\operatorname{rank}(\Id-R_j),
\]
by repeatedly using
\(\Id-AB=(\Id-A)+A(\Id-B)\).  Each reflected defect has rank \(r\);
the left side is \(\operatorname{rank}(2\Id)=n\).
\end{proof}

\begin{proposition}[Equal supported multiplicities]\label{prop:multiplicity}
Let a finite-dimensional admissible strategy and state
\(\lvert\Psi\rangle_{ABE}\) attain
\(\langle\overline{\cI}_d\rangle=M_d+1\).  Put
\[
  \rho_A=\Tr_{BE}\lvert\Psi\rangle\!\langle\Psi\rvert,\qquad
  K=\supp\rho_A,\qquad U=A_0^\dagger A_1.
\]
On \(K\), all \(d\) roots of \(z^d=(-1)^{d-1}\) occur as eigenvalues of
\(U\) with the same multiplicity.  In particular,
\[
  d\mid\dim K.
\]
No statement is made about unused sectors orthogonal to \(K\).
\end{proposition}

\begin{proof}
Equality in the two bounds
\(\langle\cI_d\rangle\leq M_d\) and
\(\langle\Reop(A_0\otimes B_d)\rangle\leq1\) is forced separately.
Every positive term in \eqref{eq:gap} therefore vanishes on the state:
\begin{align}
  P_y\lvert\Psi\rangle&=0,\label{eq:Py0}\\
  \bigl(M_d\Id-F_d(U)\bigr)\otimes\Id\lvert\Psi\rangle&=0,\label{eq:Q0}\\
  (A_0\otimes B_d\otimes\Id_E)\lvert\Psi\rangle&=\lvert\Psi\rangle.
  \label{eq:A0stab}
\end{align}
The last implication uses
\(\Id-\Reop W=\frac12(\Id-W)^\dagger(\Id-W)\) for unitary \(W\).

For any Alice operator \(R\),
\begin{equation}
  (R\otimes\Id)\lvert\Psi\rangle=0
  \quad\Longleftrightarrow\quad R|_K=0,
  \label{eq:support}
\end{equation}
as follows from a Schmidt decomposition across \(A:(BE)\).
Thus \eqref{eq:Q0} places \(K\) inside the equality spectral subspace
\[
  L=\operatorname{ran}\mathbf1_{\mathcal Z_d}(U),
  \qquad
  \mathcal Z_d=\{z:\lvert z\rvert=1,\ z^d=(-1)^{d-1}\}.
\]
Equation~\eqref{eq:A0stab} gives \(A_0K=K\).

We next justify a cancellation in \eqref{eq:Py0} without assuming that a
polar factor is globally invertible.  Let \(C_y=V_y\absop{C_y}\) be the
canonical polar decomposition, let
\[
  E_y^{\rm bad}=\mathbf1_{\{-\omega^{-y}\}}(U),\qquad
  F_y=\supp\absop{C_y^\dagger}
  =A_0(\Id-E_y^{\rm bad})A_0^\dagger.
\]
The bad phase is not in \(\mathcal Z_d\), so \(E_y^{\rm bad}K=0\).
Together with \(A_0K=K\), this gives \(F_y|_K=\Id_K\).  Also
\(F_yV_y=V_y\), because \(F_y\) is the final support of the polar partial
isometry.  Since
\[
  P_y=
  (\absop{C_y^\dagger}^{1/2}\otimes\Id)
  (\Id-V_y\otimes B_y),
\]
the vector
\[
  q_y=(\Id-V_y\otimes B_y\otimes\Id_E)\lvert\Psi\rangle
\]
belongs to \(\operatorname{ran}F_y\otimes\cH_{BE}\).
Equation~\eqref{eq:Py0} places the same vector in the kernel of
\(\absop{C_y^\dagger}^{1/2}\otimes\Id\), which is orthogonal to that range.
Hence \(q_y=0\), or
\begin{equation}
  (V_y\otimes B_y\otimes\Id_E)\lvert\Psi\rangle=\lvert\Psi\rangle.
  \label{eq:Vstab}
\end{equation}

The space \(K\) lies in the initial support of \(V_y\), so the polar factor
is an isometry there.  Tracing \eqref{eq:Vstab} over \(BE\) gives
\(\rho_A=V_y\rho_AV_y^\dagger\), and therefore \(V_yK=K\).  Hence
\[
  S_y=A_0^\dagger V_y
\]
preserves \(K\).  On \(L\), functional calculus gives
\[
  S_y=s_y(U),\qquad
  s_y(z)=\frac{1+\omega^yz}{\lvert1+\omega^yz\rvert},
  \qquad s_y(z)^2=\omega^yz.
\]
It follows that \(U=\omega^{-y}S_y^2\) preserves \(K\).  We may now let
\(E_k\) denote the spectral projections of \(U|_K\) for the \(d\) roots in
\(\mathcal Z_d\).

Put \(\eta=\e^{\pi i/d}\), and read \(y+1\) modulo \(d\), so
\(S_d=S_0\) and \(V_d=V_0\) in the next relation.  Comparing adjacent
half-angle phases shows that for each \(y\) exactly one equality root crosses
the sign cut:
\[
  \frac{1+w}{\lvert1+w\rvert}=\e^{i\theta/2}
  \quad\text{when }w=\e^{i\theta},\ -\pi<\theta<\pi.
\]
Multiplication by \(\omega\) advances the half-angle by \(\pi/d\); among the
equally spaced equality roots, exactly one principal argument wraps across
\(\pi\), contributing an additional minus sign.  Thus
\begin{equation}
  S_y^\dagger S_{y+1}
  =\eta(\Id-2E_{k(y)}),
  \qquad
  V_{y+1}=\eta V_y(\Id-2E_{k(y)}),
  \label{eq:adjacent}
\end{equation}
and \(y\mapsto k(y)\) is a permutation of the \(d\) root labels.
Raising \eqref{eq:Vstab} to the \(d\)-th power, using \(B_y^d=\Id\), and
applying \eqref{eq:support} gives \(V_y^d=\Id\) on \(K\).
Since \(\eta^d=-1\), equation~\eqref{eq:adjacent} gives
\[
  \bigl(V_y(\Id-2E_{k(y)})\bigr)^d=-\Id.
\]
Lemma~\ref{lem:reflection}, with \(n=\dim K\), yields
\[
  n\leq d\,\operatorname{rank}E_k
\]
for every \(k\).  But the \(E_k\) sum to \(\Id_K\), so their ranks sum to
\(n\).  Every inequality must therefore be an equality:
\(\operatorname{rank}E_k=n/d\) for all \(k\).
\end{proof}

\begin{remark}
Proposition~\ref{prop:multiplicity} is restricted to exact maximizers of the
\emph{augmented} functional because the proof uses the \(A_0\)-\(B_d\)
stabilizer.  Equal multiplicities do not imply a Weyl representation,
uniqueness, or self-testing; the weighted-shift family is the obstruction.
\end{remark}

\begin{thebibliography}{9}

\bibitem{PeritoEtAl2026}
I. Perito, R. D'Avino, M. Jung, P. Mironowicz, A. Ac\'in, and R. Augusiak,
\newblock Bell inequalities tailored to optimal global randomness
  certification,
\newblock \emph{arXiv:2606.21362v3} (2026),
\newblock \url{https://arxiv.org/abs/2606.21362}.

\bibitem{KriebelBound2026}
A. Kriebel,
\newblock The exact quantum value of a cyclic Bell operator,
\newblock provisional AI-assisted research note (2026),
\newblock \href{https://aleckriebel.github.io/Math/papers/cyclic-bell-tsirelson-bound/}
  {public paper page}.

\bibitem{FarkasEtAl2026}
M. Farkas, P. Mironowicz, and R. Augusiak,
\newblock Maximal global device-independent randomness from projective
  measurements in every dimension,
\newblock \emph{arXiv:2606.21369v2} (2026),
\newblock \url{https://arxiv.org/abs/2606.21369}.

\bibitem{DavinoEtAl2026}
R. D'Avino, I. Perito, P. Mironowicz, A. Ac\'in, and R. Augusiak,
\newblock Noise robustness of three outcome Bell certified quantum randomness,
\newblock \emph{arXiv:2606.21371v2} (2026),
\newblock \url{https://arxiv.org/abs/2606.21371}.

\end{thebibliography}

\end{document}
